\section{Parser Implementation}\label{sec:parser}

\subsection{Table-driven predictive parsing}

The parser is a classic non-recursive predictive parser driven by the table from
Section~\ref{sec:syntax}. It is implemented in \texttt{core/parser.py} with no generated
code and no external parsing library.

\begin{enumerate}[leftmargin=1.4em, itemsep=0.15em]
  \item Push \term{\$} and the start symbol \nonterm{S} onto the stack.
  \item Let $X$ be the stack top and $t$ the current lookahead terminal.
  \item If $X$ is a terminal it must equal $t$: pop, advance, attach the token to the
        parse tree. Otherwise \rejected\ with ``expected $X$, found $t$''.
  \item If $X$ is a nonterminal, look up $[X, t]$. Pop $X$, push the right-hand side in
        reverse, record the production id in the derivation trace. An empty cell means
        \rejected\ with the set of terminals that \emph{would} have been legal.
  \item If $X = \term{\$} = t$, \accepted.
\end{enumerate}

The parser returns the derivation as an ordered list of production identifiers
(\req{FR-25}), so every accepted statement carries a verifiable proof of \emph{why} it was
accepted, and the front end can render the parse tree from it.

\diagram{parser-automaton.png}{0.90}{Predictive parser control flow: stack, lookahead,
table lookup and the three terminating states.}

\subsection{Worked trace}

For \fw{le taximan a une voiture} the lexer produces \term{DET NOUN VERB DET NOUN} and the
parser runs as follows.

\begin{center}
\scriptsize
\begin{tabular}{@{}r l l l@{}}
\toprule
\textbf{\#} & \textbf{Stack (top first)} & \textbf{Look} & \textbf{Action} \\
\midrule
1  & \nonterm{S} \term{\$}                      & \term{DET}  & \nonterm{S} $\rightarrow$ \nonterm{Seq} \\
2  & \nonterm{Seq} \term{\$}                    & \term{DET}  & \nonterm{Seq} $\rightarrow$ \nonterm{Unit} \nonterm{Seq$'$} \\
3  & \nonterm{Unit} \nonterm{Seq$'$} \term{\$}  & \term{DET}  & \nonterm{Unit} $\rightarrow$ \nonterm{Clause} \\
4  & \nonterm{Clause} \ldots                    & \term{DET}  & \nonterm{Clause} $\rightarrow$ \nonterm{NP} \nonterm{ClauseTail} \\
5  & \nonterm{NP} \nonterm{ClauseTail} \ldots   & \term{DET}  & \nonterm{NP} $\rightarrow$ \term{DET} \nonterm{NP$'$} \\
6  & \term{DET} \nonterm{NP$'$} \ldots          & \term{DET}  & match \fw{le} \\
7  & \nonterm{NP$'$} \ldots                     & \term{NOUN} & \nonterm{NP$'$} $\rightarrow$ \term{NOUN} \\
8  & \term{NOUN} \ldots                         & \term{NOUN} & match \fw{taximan} \\
9  & \nonterm{ClauseTail} \ldots                & \term{VERB} & \nonterm{ClauseTail} $\rightarrow$ \nonterm{Predicate} \\
10 & \nonterm{Predicate} \ldots                 & \term{VERB} & \nonterm{Predicate} $\rightarrow$ \nonterm{NegOpt} \nonterm{VerbGroup} \nonterm{Comps} \\
11 & \nonterm{NegOpt} \ldots                    & \term{VERB} & \nonterm{NegOpt} $\rightarrow$ $\varepsilon$ \quad(\term{VERB} $\in$ FOLLOW) \\
12 & \nonterm{VerbGroup} \ldots                 & \term{VERB} & \nonterm{VerbGroup} $\rightarrow$ \term{VERB} \nonterm{VerbTail} \\
13 & \term{VERB} \ldots                         & \term{VERB} & match \fw{a} \\
14 & \nonterm{VerbTail} \ldots                  & \term{DET}  & \nonterm{VerbTail} $\rightarrow$ $\varepsilon$ \\
15 & \nonterm{Comps} \ldots                     & \term{DET}  & \nonterm{Comps} $\rightarrow$ \nonterm{NP} \nonterm{Comps} \\
16 & \nonterm{NP} \ldots                        & \term{DET}  & \term{DET} \nonterm{NP$'$}; match \fw{une}, \fw{voiture} \\
17 & \nonterm{Comps} \nonterm{Seq$'$} \term{\$} & \term{\$}   & \nonterm{Comps} $\rightarrow$ $\varepsilon$, \nonterm{Seq$'$} $\rightarrow$ $\varepsilon$ \\
18 & \term{\$}                                  & \term{\$}   & \accepted \\
\bottomrule
\end{tabular}
\end{center}

Steps 11 and 14 are worth pausing on: those $\varepsilon$ reductions are legal only because
\term{VERB} and \term{DET} lie in the FOLLOW sets of \nonterm{NegOpt} and
\nonterm{VerbTail} respectively. This is precisely where the FOLLOW computation earns its
place in the construction.

\subsection{Error reporting and repair suggestions}

A rejection is never a bare failure. Because the table records which terminals are legal
in each state, an empty cell yields the complete expected set, and the token's source span
gives an exact character offset (\req{FR-26}). \texttt{core/suggest.py} turns this into
concrete repairs (\req{FR-27}):

\begin{description}[leftmargin=1.4em, style=nextline, font=\ttfamily]
  \item[unknown\_word]
    The word is not in the lexicon. Near matches are computed with \texttt{difflib} over
    \emph{folded} forms at a 0.6 similarity cut-off, so an accent or spelling slip still
    finds its target.
  \item[unexpected\_token]
    The token is known but ungrammatical here. The message names the expected terminal
    classes and gives up to three real lexicon words for each, so the user is shown a
    working example rather than a category name.
  \item[trailing\_input]
    The statement parsed to completion but tokens remain --- typically two clauses run
    together with no conjunction or separator.
\end{description}

Homograph hints are emitted alongside: if \fw{la} was resolved as a determiner but \fw{là}
exists as an adverb, both readings are reported rather than one being chosen silently.

\begin{keypoint}[Suggestions are derived, not generated]
Every suggestion comes from the parser's own error report and the lexicon --- the expected
terminal set, the source span, and \texttt{difflib} over folded canonical forms. No
language model is involved, so the same input always yields the same advice
(\req{NFR-1}).
\end{keypoint}
